% !TEX program = pdflatex
\documentclass[11pt,a4paper]{article}

\usepackage[T2A]{fontenc}
\usepackage[utf8]{inputenc}
\usepackage[russian]{babel}
\usepackage{amsmath,amssymb}
\usepackage{array,booktabs,longtable,tabularx}
\usepackage{geometry}
\usepackage{microtype}
\usepackage{xcolor}
\usepackage[hidelinks]{hyperref}

\geometry{margin=18mm}
\setlength{\parindent}{1.25em}
\setlength{\parskip}{0.25em}
\renewcommand{\arraystretch}{1.12}
\newcolumntype{Y}{>{\raggedright\arraybackslash}X}
\newcommand{\code}[1]{\texttt{#1}}
\newcommand{\qsi}{Q_{\mathrm{SI}}}
\newcommand{\qmi}{Q_{\mathrm{MI}}}
\newcommand{\snf}{S/N/F}

\title{Отчёт по численным экспериментам ADP\\
\large Single-index и multi-index модели}
\author{}
\date{2 сентября 2026 г.}

\begin{document}
\maketitle

\begin{abstract}
В отчёте объединены результаты 62 серий из шести переданных наборов
артефактов. Для каждой группы приведены постановка эксперимента,
фактически использованные входные параметры и агрегированные результаты.
Качество оценки, техническая сходимость и численные ошибки рассматриваются
раздельно. Всего обработано 3498 запусков: 847 завершились со статусом
\code{success}, 2648 --- со статусом \code{nonconverged}, три --- со
статусом \code{numerical\_failure}.
\end{abstract}

\section{Постановка и метрики}

Для single-index модели генерировались данные
\begin{equation}
    Y_i = f(\beta_*^\top X_i) + \varepsilon_i,
    \qquad \|\beta_*\|_2=1,
\end{equation}
а для multi-index модели размерности $m=2$
\begin{equation}
    Y_i = f(B_*^\top X_i) + \varepsilon_i,
    \qquad B_*^\top B_*=I_2.
\end{equation}
Сигнал перед добавлением шума стандартизировался. При заданном $\tau$
признаки строились по формуле
\begin{equation}
    X_i = \sigma_X\{\tau z_{0i}+(1-\tau)z_i\},
\end{equation}
где $z_{0i}$ --- общий гауссовский фактор, а координаты $z_i$ независимы.
Шум во включённых сериях гауссовский. Основные функции связи:
\begin{align}
 f_{\sin}(x;s)&=\sin(sx), &
 f_{x\sin}(x;s)&=x\sin(sx),\\
 f_{\mathrm{add}}(z_1,z_2;s)&=z_1^2+\sin(sz_2), &
 f_{\mathrm{mult}}(z_1,z_2;s)&=z_1\sin(sz_2).
\end{align}

Качество single-index оценки измерялось величиной
\begin{equation}
    \qsi=|\beta_*^\top\widehat\beta|,
\end{equation}
где большее значение лучше. Для multi-index использовалось RMS-расстояние
по главным углам
\begin{equation}
    \qmi=
    \left(\frac{1}{m}\sum_{j=1}^{m}\sin^2\theta_j\right)^{1/2},
\end{equation}
где меньшее значение лучше. В сериях с формальным порогом качественным
считалось $\qsi\geq0.9$ или $\qmi\leq0.1$. Обозначение $\snf$ в таблицах
означает число запусков со статусами \code{success},
\code{nonconverged} и \code{numerical\_failure}. Прохождение порога качества
не означает технической сходимости; восстановление требует обоих условий.


\section{Функции связи и уточняющие multi-index серии}

\subsection{Входные параметры}

Single-index серии по функциям использовали $n=600$, $d=30$, $\tau=0.4$,
$\sigma_\varepsilon=0.2$, $N_{\mathrm{loc}}=20$,
$N_{\mathrm{lin}}=90$, $N_J=200$, $N_\phi=20$, пять внешних и восемь
внутренних шагов. Multi-index серии использовали $n=1000$, базово $d=25$,
$\tau=0.4$, $\sigma_\varepsilon=0.2$, $N_{\mathrm{loc}}=12$,
$N_{\mathrm{lin}}=85$, $N_J=500$, $N_\phi=30$, три внешних и пять
внутренних шагов. Для всех точек выполнено по три повтора.

\begin{longtable}{p{0.31\textwidth}p{0.18\textwidth}p{0.43\textwidth}}
\caption{Варьируемые входы функциональных и multi-index серий.}
\label{tab:function-inputs}\\
\toprule
Серия & Фактор & Уровни \\
\midrule
\endfirsthead
\toprule
Серия & Фактор & Уровни \\
\midrule
\endhead
\code{si-function-classes} & $f$ & linear, quadratic, square, cubic, quartic, $\sin$, $\cos$, $x\sin$, $\tanh$, $|x|$, ReLU, Gaussian bump \\
\code{si-parity-frequency} & $f,s$ & $\sin(sx)$ и $\cos(sx)$; $s=0.5,1,1.5,2,2.5,3,3.5$ \\
\code{si-gradient-support} & $f,s$ & $\tanh(sx)$ и $\exp(-(sx)^2/2)$; $s=0.25,0.5,0.75,1,1.5,2,3$ \\
\code{si-init-by-function} & $f$, init & шесть функций; local и random \\
\code{mi-focus-d} & $d$ & 20,22,24,25,26,28,30; $N_{\mathrm{lin}}=d+60$ \\
\code{mi-focus-frequency-additive} & $s$ & 2,2.15,2.3,2.5,2.7,2.85,3 \\
\code{mi-focus-noise} & $\sigma_\varepsilon$ & 0.6,0.65,0.7,0.75,0.8,0.9,1 \\
\code{mi-focus-nphi} & $N_\phi$ & 20,24,28,30,32,36,40 \\
\code{mi-focus-init} & init & local, random \\
\code{mi-focus-tensor} & tensor & full, orthogonal \\
\bottomrule
\end{longtable}

\subsection{Результаты}

\begin{table}[ht]
\centering
\small
\begin{tabular}{l|l|r|r|r}
\hline
Multi-index серия & Лучший уровень & $\qmi$ & Порог & $\snf$ \\
\hline
\code{mi-focus-d} & $d=22$ & 0.0955 & 4/21 & 0/21/0 \\
\code{mi-focus-frequency-additive} & $s=2.15$ & 0.0628 & 15/21 & 0/21/0 \\
\code{mi-focus-noise} & $\sigma_\varepsilon=0.6$ & 0.2115 & 0/21 & 0/21/0 \\
\code{mi-focus-nphi} & $N_\phi=28$ & 0.0728 & 9/21 & 0/21/0 \\
\code{mi-focus-init} & local & 0.0921 & --- & 0/6/0 \\
\code{mi-focus-tensor} & full & 0.1408 & --- & 0/6/0 \\
\hline
\end{tabular}
\caption{Итог multi-index уточнений. В столбце ``Порог'' указано число
точек с $\qmi\leq0.1$.}
\label{tab:function-mi-results}
\end{table}

\begin{table}[ht]
\centering
\small
\begin{tabularx}{\textwidth}{Yrrrr}
\toprule
Single-index серия & Запусков & $\qsi$ (медиана) & Порог & $\snf$ \\
\midrule
\code{si-function-classes} & 36 & 0.9984 & не задан & 5/31/0 \\
\code{si-gradient-support} & 42 & 0.9985 & 42/42 & 8/34/0 \\
\code{si-init-by-function} & 36 & 0.9984 & не задан & 11/25/0 \\
\code{si-parity-frequency} & 42 & 0.9990 & 42/42 & 5/37/0 \\
\bottomrule
\end{tabularx}
\caption{Итог single-index серий по функциям связи.}
\label{tab:function-si-results}
\end{table}

Во всех 12 классах функций медиана $\qsi$ лежит между 0.9923 и 0.9988.
В сериях чётности/частоты и ширины информативной производной все 84 запуска
прошли порог качества. Предсказанного ухудшения при росте частоты до 3.5 и
сужении области производной до масштаба 3 на этой сетке не обнаружено.
Локальная инициализация заметно помогла для Gaussian bump
($0.9993$ против $0.9789$), но для остальных проверенных функций разность
медиан мала.

Для multi-index зависимость более выражена. При $d=20,22$ порог прошли по два
запуска из трёх, а при $d\geq24$ медиана уже выше 0.1. Для additive-link
лучший уровень $s=2.15$ дал медиану 0.0628. При шуме
$\sigma_\varepsilon\geq0.6$ ни один запуск не прошёл порог. По числу
направлений наблюдается локальный минимум около $N_\phi=28$--32, но все
multi-index запуски остались технически несошедшимися.

\section{Повторные single-index границы: \code{single 2}}

\subsection{Входные параметры}

Базовая конфигурация: $n=1000$, $d=100$, $\sigma_X=1$, $\tau=0.4$,
$\sigma_\varepsilon=0.2$, $f(x)=x\sin(3x)$,
$N_{\mathrm{loc}}=20$, $N_{\mathrm{lin}}=160$, $N_J=500$,
$N_\phi=20$, три внешних и три внутренних шага, local initialization и
localized directions. 
\begin{table}[ht]
\centering
\small
\begin{tabular}{l|l|l}
\hline
Серия & Фактор & Уровни \\
\hline
\code{si-focus-d} & $d$ & 40,44,48,50,52,56,60 \\
\code{si-focus-frequency} & $f,s$ & $\sin(sx)$, $x\sin(sx)$; $s=2.5,2.75,3,3.25,3.5,3.75,4$ \\
\code{si-focus-tau} & $\tau$ & 0.6,0.65,0.7,0.75,0.8,0.85,0.9 \\
\code{si-focus-displacement} & смещение & 0.5,0.6,0.7,0.75,0.8,0.9,1 \\
\hline
\end{tabular}
\caption{Входные сетки повторных single-index серий.}
\label{tab:single2-inputs}
\end{table}

\subsection{Результаты}

\begin{table}[ht]
\centering
\small
\begin{tabular}{c|rrrrrrr}
\hline
$d$ & 40 & 44 & 48 & 50 & 52 & 56 & 60 \\
\hline
медиана $\qsi$ & .604 & .743 & .433 & .367 & .378 & .856 & .396 \\
порог / 3 & 0 & 0 & 0 & 0 & 0 & 1 & 0 \\
время, с & 1.16 & 1.29 & 1.37 & 1.46 & 1.47 & 1.72 & 1.98 \\
\hline
\end{tabular}
\caption{Повтор границы по размерности. Все 21 запуска имеют статус
\code{nonconverged}.}
\label{tab:single2-d}
\end{table}

\begin{table}[ht]
\centering
\small
\begin{tabular}{c|rrrrrrr}
\hline
$\tau$ & .60 & .65 & .70 & .75 & .80 & .85 & .90 \\
\hline
медиана $\qsi$ & .066 & .480 & .934 & .958 & .975 & .971 & .973 \\
порог / 3 & 0 & 0 & 3 & 3 & 3 & 3 & 3 \\
\hline
\end{tabular}
\caption{Повтор границы по общему фактору. Все 21 запуска имеют статус
\code{nonconverged}.}
\label{tab:single2-tau}
\end{table}

\begin{table}[ht]
\centering
\small
\begin{tabular}{l|rrrrrrr}
\hline
$s$ & 2.50 & 2.75 & 3.00 & 3.25 & 3.50 & 3.75 & 4.00 \\
\hline
$\sin(sx)$: медиана $\qsi$ & .967 & .941 & .899 & .801 & .452 & .258 & .124 \\
$\sin(sx)$: порог / 3 & 3 & 3 & 1 & 0 & 0 & 0 & 0 \\
$x\sin(sx)$: медиана $\qsi$ & .440 & .654 & .735 & .757 & .752 & .708 & .535 \\
$x\sin(sx)$: порог / 3 & 0 & 0 & 0 & 1 & 1 & 0 & 0 \\
\hline
\end{tabular}
\caption{Повтор частотной границы; все 42 запуска технически не сошлись.}
\label{tab:single2-frequency}
\end{table}

\begin{table}[ht]
\centering
\small
\begin{tabular}{c|rrrrrrr}
\hline
Смещение & .50 & .60 & .70 & .75 & .80 & .90 & 1.00 \\
\hline
медиана $\qsi$ & .064 & .053 & .059 & .036 & .064 & .055 & .151 \\
порог / доступные & 0/3 & 0/3 & 0/3 & 0/3 & 0/3 & 0/3 & 0/2 \\
$\snf$ & 0/3/0 & 0/3/0 & 0/3/0 & 0/3/0 & 0/3/0 & 0/3/0 & 0/2/1 \\
\hline
\end{tabular}
\caption{Повтор границы смещения центров.}
\label{tab:single2-displacement}
\end{table}

Повтор подтвердил резкую границу для $\sin(sx)$ между $s=2.75$ и $s=3.25$
и переход по общему фактору между $\tau=0.65$ и $\tau=0.70$. Граница по
размерности не воспроизвелась монотонно: только один из 21 запусков прошёл
порог. Смещение центров не улучшило качество; при смещении 1.0 один из трёх
запусков завершился численной ошибкой компактной поддержки. Из-за трёх вместо
пяти повторов эти выводы остаются предварительными.

\section{Масштабирование \texorpdfstring{$N_\phi$, $N_{\mathrm{loc}}$ и
$N_J$}{Nphi, Nloc и NJ}}

\subsection{Описание и входные параметры}

Для single-index и multi-index использован полный факторный дизайн
\begin{equation}
    n\in\{500,1000,2000\},\qquad
    d\in\{10,25,50\},
\end{equation}
с пятью повторами на каждую точку. Варьировался ровно один из параметров:
\begin{align}
 N_\phi &\in \{5,10,20,30,40,60\} &&\text{для SI},\\
 N_\phi &\in \{3,5,10,20,30,40\} &&\text{для MI},\\
 N_{\mathrm{loc}} &\in \{7,10,15,20,30,45\},\\
 N_J &= \left\lceil c\,n/N_{\mathrm{loc}}\right\rceil,
 &c&\in\{1,2,4,8,12\}.
\end{align}
Вне изменяемого фактора использовались $N_{\mathrm{loc}}=15$,
$N_{\mathrm{lin}}=d+60$, $N_J=\lceil n/5\rceil$, $N_\phi=20$,
$\sigma_X=1$, $\tau=0$, $\sigma_\varepsilon=0.2$, пять внешних и восемь
внутренних шагов. SI использовал $\sin(x)$, local initialization и localized
directions; MI --- additive-link, pilot initialization и isotropic directions.

\begin{table}[ht]
\centering
\small
\begin{tabular}{l|l|l}
\hline
Модель & Selector & Изменяемый параметр \\
\hline
SI & \code{scale-si-nphi} & $N_\phi$ \\
SI & \code{scale-si-nloc} & $N_{\mathrm{loc}}$ \\
SI & \code{scale-si-nj} & $N_J=\lceil c n/N_{\mathrm{loc}}\rceil$ \\
MI & \code{scale-mi-nphi} & $N_\phi$ \\
MI & \code{scale-mi-nloc} & $N_{\mathrm{loc}}$ \\
MI & \code{scale-mi-nj} & $N_J=\lceil c n/N_{\mathrm{loc}}\rceil$ \\
\hline
\end{tabular}
\caption{Параметрические серии и изменяемые входы.}
\label{tab:param-selectors}
\end{table}

\subsection{Результаты}

\begin{longtable}{llrrrrrr}
\caption{Параметрические результаты, агрегированные по девяти парам $(n,d)$.}
\label{tab:param-results}\\
\toprule
Модель & Уровень & Медиана $Q$ & Порог & $S$ & $N$ & Время, с & Память, MiB \\
\midrule
\endfirsthead
\toprule
Модель & Уровень & Медиана $Q$ & Порог & $S$ & $N$ & Время, с & Память, MiB \\
\midrule
\endhead
MI & $N_\phi=3$  & .7249 & 8/45  & 0 & 45 & 3.217 & 7.76 \\
MI & $N_\phi=5$  & .1877 & 17/45 & 0 & 45 & 2.357 & 7.76 \\
MI & $N_\phi=10$ & .0984 & 23/45 & 0 & 45 & 2.336 & 9.34 \\
MI & $N_\phi=20$ & .0851 & 27/45 & 0 & 45 & 2.379 & 17.15 \\
MI & $N_\phi=30$ & .0788 & 28/45 & 0 & 45 & 3.369 & 24.96 \\
MI & $N_\phi=40$ & .0785 & 29/45 & 0 & 45 & 3.650 & 30.66 \\
\midrule
MI & $N_{\mathrm{loc}}=7$  & .1165 & 19/45 & 0 & 45 & 2.882 & 9.32 \\
MI & $N_{\mathrm{loc}}=10$ & .0855 & 29/45 & 0 & 45 & 2.561 & 17.15 \\
MI & $N_{\mathrm{loc}}=15$ & .0801 & 29/45 & 0 & 45 & 2.705 & 17.15 \\
MI & $N_{\mathrm{loc}}=20$ & .0775 & 31/45 & 0 & 45 & 2.591 & 17.15 \\
MI & $N_{\mathrm{loc}}=30$ & .0724 & 32/45 & 0 & 45 & 2.789 & 17.15 \\
MI & $N_{\mathrm{loc}}=45$ & .0706 & 34/45 & 0 & 45 & 3.018 & 17.15 \\
\midrule
MI & $c=1$  & .0914 & 25/45 & 0 & 45 & 2.460 & 10.36 \\
MI & $c=2$  & .0830 & 30/45 & 0 & 45 & 2.122 & 16.86 \\
MI & $c=4$  & .0685 & 31/45 & 0 & 45 & 3.343 & 17.44 \\
MI & $c=8$  & .0660 & 29/45 & 0 & 45 & 5.446 & 20.51 \\
MI & $c=12$ & .0786 & 28/45 & 0 & 45 & 6.555 & 30.66 \\
\midrule
SI & $N_\phi=5$  & .9991 & 45/45 & 0  & 45 & .975 & 9.16 \\
SI & $N_\phi=10$ & .9992 & 45/45 & 3  & 42 & .998 & 9.60 \\
SI & $N_\phi=20$ & .9993 & 45/45 & 12 & 33 & 1.109 & 17.40 \\
SI & $N_\phi=30$ & .9993 & 45/45 & 19 & 26 & 1.232 & 25.21 \\
SI & $N_\phi=40$ & .9993 & 45/45 & 18 & 27 & 1.346 & 33.02 \\
SI & $N_\phi=60$ & .9994 & 45/45 & 27 & 18 & 2.021 & 36.63 \\
\midrule
SI & $N_{\mathrm{loc}}=7$  & .9989 & 45/45 & 0  & 45 & .891 & 9.44 \\
SI & $N_{\mathrm{loc}}=10$ & .9991 & 45/45 & 0  & 45 & .936 & 9.44 \\
SI & $N_{\mathrm{loc}}=15$ & .9991 & 45/45 & 9  & 36 & 1.026 & 17.40 \\
SI & $N_{\mathrm{loc}}=20$ & .9992 & 45/45 & 25 & 20 & 1.113 & 17.40 \\
SI & $N_{\mathrm{loc}}=30$ & .9993 & 45/45 & 34 & 11 & 1.138 & 17.40 \\
SI & $N_{\mathrm{loc}}=45$ & .9993 & 45/45 & 37 & 8  & 1.224 & 17.40 \\
\midrule
SI & $c=1$  & .9990 & 45/45 & 5  & 40 & .679 & 12.27 \\
SI & $c=2$  & .9992 & 45/45 & 12 & 33 & .910 & 17.12 \\
SI & $c=4$  & .9992 & 45/45 & 13 & 32 & 1.189 & 17.69 \\
SI & $c=8$  & .9993 & 45/45 & 15 & 30 & 2.166 & 24.45 \\
SI & $c=12$ & .9993 & 45/45 & 20 & 25 & 2.857 & 36.63 \\
\bottomrule
\end{longtable}

Для SI качество практически насыщено на всей сетке: все 765 запусков прошли
порог $\qsi\geq0.9$. Рост $N_{\mathrm{loc}}$ и $N_\phi$ повышает долю
технически сошедшихся запусков, но увеличивает время; ни один уровень не дал
полной сходимости. Для MI все 765 запусков остались \code{nonconverged}.
Наиболее выражен эффект $N_\phi$: переход от 3 к 20 уменьшает медиану
$\qmi$ с 0.7249 до 0.0851, после чего выигрыш мал относительно роста памяти.
Для $N_{\mathrm{loc}}$ качество улучшается плавно вплоть до 45. Для числа
центров $c=4$ даёт почти то же качество, что $c=8$, при заметно меньших
времени и памяти; $c=12$ уже ухудшает медиану. Эти уровни являются
инженерными компромиссами screening, а не подтверждёнными оптимумами, потому
что техническая сходимость MI не достигнута.

\section{Подробная multi-index граница по
\texorpdfstring{$n$ и $d$}{n и d}}

\subsection{Входные параметры}

В серии \code{mi-boundary-nd} estimator фиксирован, а меняются только размер
выборки и ambient dimension. Фактические параметры приведены в
таблице~\ref{tab:mi-nd-inputs}. Выполнено 315 запусков: 63 пары $(n,d)$ по
пять seed

\begin{table}[ht]
\centering
\small
\begin{tabular}{l|l}
\hline
Параметр & Значение \\
\hline
$n$ & 400,600,800,1000,1400,1800,2400 \\
$d$ & 10,15,20,25,30,35,40,50,60 \\
$m$ & 2 \\
$\sigma_X,\tau,\sigma_\varepsilon$ & $1,0,0.2$ \\
Функция связи & $z_1^2+\sin(z_2)$, затем стандартизация \\
$N_{\mathrm{loc}}$ & 15 \\
$N_{\mathrm{lin}}$ & $d+60$ \\
$N_J$ & $\lceil n/5\rceil$ \\
$N_\phi$ & 20 \\
Внешних / внутренних шагов & 5 / 8 \\
$\lambda,a,h_{\min}$ & $0.05,\sqrt2,3/\sqrt n$ \\
Инициализация / направления & pilot / isotropic \\
Тензор / выбор шага & orthogonal / best \\
Повторов & 5 фактически, 10 в \code{full\_runs} \\
\hline
\end{tabular}
\caption{Входные параметры \code{mi-boundary-nd}.}
\label{tab:mi-nd-inputs}
\end{table}

\subsection{Результаты}

В таблице~\ref{tab:mi-nd-matrix} каждая ячейка содержит медиану $\qmi$ и в
скобках число запусков из пяти, прошедших порог $\qmi\leq0.1$.

\begin{table}[ht]
\centering
\small
\begin{tabular}{c|rrrrrrr}
\hline
$d\backslash n$ & 400 & 600 & 800 & 1000 & 1400 & 1800 & 2400 \\
\hline
10 & .059(5) & .053(5) & .044(5) & .057(5) & .039(5) & .030(5) & .032(5) \\
15 & .089(3) & .055(5) & .057(5) & .058(5) & .049(5) & .047(5) & .037(5) \\
20 & .106(1) & .073(4) & .084(3) & .079(3) & .062(4) & .050(5) & .049(5) \\
25 & .111(1) & .108(1) & .082(4) & .083(3) & .070(5) & .054(5) & .057(4) \\
30 & .167(0) & .124(0) & .088(3) & .093(4) & .079(4) & .072(4) & .071(4) \\
35 & .200(0) & .153(0) & .111(2) & .102(2) & .097(3) & .079(4) & .068(4) \\
40 & .212(0) & .143(0) & .117(0) & .097(3) & .088(4) & .081(4) & .079(5) \\
50 & .748(0) & .226(0) & .148(0) & .121(1) & .193(1) & .095(3) & .133(1) \\
60 & .441(0) & .204(0) & .175(0) & .156(0) & .127(1) & .101(2) & .091(4) \\
\hline
\end{tabular}
\caption{Медианная ошибка подпространства и число прохождений порога для
всех 63 пар $(n,d)$.}
\label{tab:mi-nd-matrix}
\end{table}

\begin{table}[ht]
\centering
\small
\begin{tabular}{r|rrrr}
\hline
$n$ & медиана $\qmi$ & Порог & Время, с & Память, MiB \\
\hline
400  & .1511 & 10/45 & 2.598 & 7.10 \\
600  & .1219 & 15/45 & 2.505 & 10.52 \\
800  & .1070 & 22/45 & 2.401 & 13.95 \\
1000 & .0972 & 26/45 & 2.790 & 17.37 \\
1400 & .0809 & 32/45 & 2.822 & 24.22 \\
1800 & .0722 & 37/45 & 2.969 & 30.90 \\
2400 & .0625 & 37/45 & 3.794 & 44.08 \\
\hline
\end{tabular}
\qquad
\begin{tabular}{r|rrrr}
\hline
$d$ & медиана $\qmi$ & Порог & Время, с & Память, MiB \\
\hline
10 & .0436 & 35/35 & 1.830 & 13.14 \\
15 & .0580 & 33/35 & 1.998 & 15.08 \\
20 & .0693 & 25/35 & 2.341 & 15.08 \\
25 & .0819 & 23/35 & 2.594 & 17.15 \\
30 & .0926 & 19/35 & 2.826 & 17.37 \\
35 & .1111 & 15/35 & 3.087 & 17.60 \\
40 & .1122 & 16/35 & 3.316 & 17.83 \\
50 & .1570 & 6/35  & 4.213 & 18.28 \\
60 & .1790 & 7/35  & 4.517 & 18.73 \\
\hline
\end{tabular}
\caption{Маргинальные результаты по $n$ и $d$.}
\label{tab:mi-nd-marginals}
\end{table}

\begin{table}[ht]
\centering
\small
\begin{tabular}{r|rrrrrrrrr}
\hline
$d$ & 10 & 15 & 20 & 25 & 30 & 35 & 40 & 50 & 60 \\
\hline
Минимальное $n$ & 400 & 400 & 600 & 800 & 800 & 1400 & 1000 & 1800 & 2400 \\
Порог / 5 & 5 & 3 & 4 & 4 & 3 & 3 & 3 & 3 & 4 \\
\hline
\end{tabular}
\caption{Первое $n$ в проверенной сетке, при котором медиана
$\qmi\leq0.1$. Из-за пяти seed граница немонотонна.}
\label{tab:mi-nd-boundary}
\end{table}

Рост $n$ улучшает оценку: маргинальная медиана снижается с 0.1511 до 0.0625,
тогда как traced peak memory возрастает с 7.10 до 44.08 MiB. Рост $d$ ухудшает
качество и увеличивает время: медиана меняется от 0.0436 при $d=10$ до
0.1790 при $d=60$, а время --- от 1.83 до 4.52 с. При одинаковом отношении
$n/d=40$ медианы остаются ниже 0.1, но меняются примерно от 0.055 до 0.097;
следовательно, одного отношения $n/d$ недостаточно для точного описания
конечной выборки. Наблюдаются отдельные немонотонности, особенно при $d=50$,
поэтому таблицу~\ref{tab:mi-nd-boundary} следует считать эмпирической
границей screening, а не строгим законом.


\clearpage
\section{Зависимость \texorpdfstring{$N_{\mathrm{loc}}$}{Nloc} от
\texorpdfstring{$n$ и $d$}{n и d} в single-index модели}

\subsection{Описание и входные параметры}

Серия \code{si-nloc-nd} проверяет гипотезу о том, что для устойчивой
технической сходимости при росте размерности требуется больше локальных
точек $N_{\mathrm{loc}}$, а после насыщения качества дальнейшее увеличение
$N_{\mathrm{loc}}$ в основном повышает вычислительную стоимость. В каждой
точке сетки оценивалась одна и та же single-index модель; менялись $n$, $d$
и $N_{\mathrm{loc}}$. Выполнено 1050 запусков: 210 комбинаций параметров по
пять seed.

\begin{table}[ht]
\centering
\small
\begin{tabular}{l|l}
\hline
Параметр & Значение \\
\hline
$n$ & 400,700,1000,1500,2200 \\
$d$ & 10,20,30,40,50,60 \\
$N_{\mathrm{loc}}$ & 10,15,20,25,30,40,50 \\
$m$ & 1 \\
$\sigma_X,\tau,\sigma_\varepsilon$ & $1,0,0.2$ \\
Функция связи & $\sin(x)$, затем стандартизация \\
$N_{\mathrm{lin}}$ & $d+60$ \\
$N_J$ & $\lceil n/5\rceil$ \\
$N_\phi$ & 20 \\
Внешних / внутренних шагов & 5 / 8 \\
$\lambda,a,h_{\min}$ & $0.05,\sqrt2,3/\sqrt n$ \\
Инициализация / направления & local / localized \\
Выбор шага & best \\
Решатель & LSMR, tolerance $10^{-6}$ \\
Локальный ridge / batch size & $10^{-8}$ / 32 \\
Ядро & Epanechnikov \\
Повторов & 5 фактически, 10 в \code{full\_runs} \\
\hline
\end{tabular}
\caption{Входные параметры \code{si-nloc-nd}.}
\label{tab:si-nloc-nd-inputs}
\end{table}

\subsection{Результаты}

Во всех 1050 запусках выполнен порог $\qsi\geq0.9$; численных отказов не
было. Статус \code{success} получили 569 запусков, а
\code{nonconverged} --- 481. Поэтому число восстановлений, одновременно
удовлетворяющих критериям качества и технической сходимости, равно 569.

\begin{longtable}{c r r r r r r}
\caption{Маргинальные результаты \code{si-nloc-nd}. $S$ и $N$ --- числа
статусов \code{success} и \code{nonconverged}.}
\label{tab:si-nloc-nd-marginals}\\
\toprule
Фактор & Уровень & Запуски & $S$ & $N$ & Медиана $\qsi$ & Время, с / MiB \\
\midrule
\endfirsthead
\toprule
Фактор & Уровень & Запуски & $S$ & $N$ & Медиана $\qsi$ & Время, с / MiB \\
\midrule
\endhead
$n$ & 400  & 210 & 54  & 156 & .9974 & .999 / 7.31 \\
    & 700  & 210 & 104 & 106 & .9985 & 1.090 / 12.58 \\
    & 1000 & 210 & 124 & 86  & .9988 & 1.240 / 17.75 \\
    & 1500 & 210 & 149 & 61  & .9993 & 1.637 / 26.31 \\
    & 2200 & 210 & 138 & 72  & .9995 & 2.469 / 44.32 \\
\midrule
$d$ & 10 & 175 & 129 & 46  & .9997 & .716 / 16.73 \\
    & 20 & 175 & 122 & 53  & .9995 & .968 / 17.18 \\
    & 30 & 175 & 98  & 77  & .9991 & 1.190 / 17.63 \\
    & 40 & 175 & 90  & 85  & .9987 & 1.410 / 18.08 \\
    & 50 & 175 & 69  & 106 & .9986 & 1.633 / 18.53 \\
    & 60 & 175 & 61  & 114 & .9982 & 1.722 / 18.98 \\
\midrule
$N_{\mathrm{loc}}$ & 10 & 150 & 5   & 145 & .9989 & 1.086 / 17.75 \\
    & 15 & 150 & 29  & 121 & .9989 & 1.217 / 17.85 \\
    & 20 & 150 & 74  & 76  & .9990 & 1.285 / 17.85 \\
    & 25 & 150 & 98  & 52  & .9991 & 1.344 / 17.85 \\
    & 30 & 150 & 113 & 37  & .9991 & 1.393 / 17.85 \\
    & 40 & 150 & 122 & 28  & .9991 & 1.412 / 17.85 \\
    & 50 & 150 & 128 & 22  & .9992 & 1.443 / 17.85 \\
\bottomrule
\end{longtable}

\begin{table}[ht]
\centering
\small
\begin{tabular}{r|rrrrrrr}
\hline
$d\backslash N_{\mathrm{loc}}$ & 10 & 15 & 20 & 25 & 30 & 40 & 50 \\
\hline
10 & 0 & 10 & 21 & 24 & 24 & 25 & 25 \\
20 & 5 & 10 & 16 & 22 & 21 & 23 & 25 \\
30 & 0 & 4  & 12 & 16 & 20 & 23 & 23 \\
40 & 0 & 5  & 13 & 16 & 17 & 19 & 20 \\
50 & 0 & 0  & 7  & 9  & 17 & 18 & 18 \\
60 & 0 & 0  & 5  & 11 & 14 & 14 & 17 \\
\hline
\end{tabular}
\caption{Число статусов \code{success} из 25 запусков для каждой пары
$(d,N_{\mathrm{loc}})$; объединены пять уровней $n$.}
\label{tab:si-nloc-d-matrix}
\end{table}

\begin{table}[ht]
\centering
\small
\begin{tabular}{r|rrrrrrr}
\hline
$n\backslash N_{\mathrm{loc}}$ & 10 & 15 & 20 & 25 & 30 & 40 & 50 \\
\hline
400  & 0 & 1 & 6  & 9  & 11 & 14 & 13 \\
700  & 0 & 4 & 12 & 18 & 21 & 24 & 25 \\
1000 & 1 & 7 & 16 & 20 & 25 & 25 & 30 \\
1500 & 2 & 8 & 22 & 27 & 30 & 30 & 30 \\
2200 & 2 & 9 & 18 & 24 & 26 & 29 & 30 \\
\hline
\end{tabular}
\caption{Число статусов \code{success} из 30 запусков для каждой пары
$(n,N_{\mathrm{loc}})$; объединены шесть уровней $d$.}
\label{tab:si-nloc-n-matrix}
\end{table}

Качество оценки насыщено на всей сетке: маргинальные медианы $\qsi$
находятся в диапазоне от 0.9974 до 0.9997. При увеличении
$N_{\mathrm{loc}}$ число технически сошедшихся запусков возрастает с 5 из
150 до 128 из 150, медианное время --- с 1.086 до 1.443 с, а медианная
traced peak memory остаётся практически неизменной: 17.75--17.85 MiB.
Таким образом, в проверенной сетке рост $N_{\mathrm{loc}}$ влияет прежде
всего на техническую сходимость, а не на уже насыщенное качество.

Рост $d$ сопровождается снижением числа \code{success} со 129 до 61 из 175
и ростом медианного времени с 0.716 до 1.722 с. Увеличение $n$ в целом
улучшает сходимость до $n=1500$, но при $n=2200$ число успешных запусков
снижается со 149 до 138; следовательно, зависимость не является строго
монотонной. Одновременно медианная traced peak memory растёт с 7.31 до
44.32 MiB. Матрица~\ref{tab:si-nloc-d-matrix} качественно поддерживает
гипотезу о необходимости большего $N_{\mathrm{loc}}$ при росте $d$, однако
даже $N_{\mathrm{loc}}=50$ не обеспечивает полной сходимости. Поскольку
выполнено пять повторов вместо запланированных десяти, этот вывод следует
считать предварительным.



\end{document}
